\subsection{Sixteen actual-source candidates in the retained mixed quotient}

\begin{statement}{D23.C25}{Proved}{The normalized degree-72 covariance gate}
\end{statement}
Use the genuine degree-1344 Hurwitz action of D23.P17--P18, with
left permutation composition, and put
\[
 X=C^2,\qquad Y=(ABA)^2,\qquad R=f(X,Y),\qquad J=\langle X,Y\rangle.
\]
Suppose an actual multiplier-one full-GT element has
$(\rho,a_3,a_2)=(0,\epsilon,0)$, where $\epsilon\in\{0,1\}$.
Let $V$ be its actual permutation supplied by P18.  The two
degree-72 $J$-orbits numbered 9 and 10 are each preserved by $V$.
On either orbit, exactly sixteen values of $R$ satisfy the following
necessary conditions together:
\begin{enumerate}
\item $R$ belongs to the derived subgroup of the marked mixed image;
\item its marked $(p,t)=((123),(1234))$ quotient is the identity;
\item there is a permutation $T$ fixing the 24 points in orthogonal
      blocks $7,8,10$ and satisfying
      $TXT^{-1}=RXR^{-1}$ and $TYT^{-1}=Y$.
\end{enumerate}
For each such $R$, the permutation $T$ is unique.  The sixteen
values form an elementary abelian group of order sixteen.  Their
four distinct $T$ values have orders $1,2,2,2$, with four values
of $R$ over each $T$.

The two marked degree-72 representations have the same kernel.
Their sixteen surviving values therefore match in a graph with
sixteen pairs.  In the saved lexicographic ordering of their four
$T$ values, the only possible pairs are
\[
                   (0,0),\quad(1,3),\quad(2,1),\quad(3,2).
\]
All these conclusions are identical for $\epsilon=0$ and
$\epsilon=1$.  They impose necessary finite restrictions; they do
not construct an actual associator or prove the existence of the
normalized unequal-bit state.

\begin{proof}
The reversed standard free-group automorphism is
\[
            x\longmapsto f(x,y)x f(x,y)^{-1},
            \qquad y\longmapsto y.
\]
Indeed it is obtained from the usual gentle automorphism by
interchanging the two free generators and using the two-cycle.
It is an automorphism of the free profinite group.  Its images in
any finite quotient therefore still generate that quotient.
P18 and the fixed three-strand full twist give
\[
                VXV^{-1}=RXR^{-1},\qquad VYV^{-1}=Y.
\]
Consequently $VJV^{-1}=J$, so $V$ permutes the $J$-orbits.

The exact orbit enumeration has two orbits of degree 72.  Each
meets each of blocks $1,2,3,4,5,6,7,8,10$ in eight points and
meets no other block.  P21 gives the actual maps on blocks
$7,8,10,11$.  Substitution of those literal permutations, for
either cubic bit, fixes every one of the 24 indicated points in
each degree-72 orbit.  A normalizer of $J$ sends a $J$-orbit onto
a $J$-orbit; since these orbits contain fixed points, each is
preserved individually.

On each component, C21 gives the mixed image order 995328 and
derived order 82944, and the exact marked $S_4$ quotient.
The supplied FP.37.2 formula for full GT is
$f(p,t)=t^{2a_2}p^\rho$.  Thus the normalized state has identity
image in that quotient.  The joined permutation representation
on $72+4$ points has the same group order.  Enumeration of its
derived subgroup finds exactly 6912 values with last four
coordinates fixed.

For each of these 6912 values $R$, solve the two displayed
intertwining equations by propagation.  Begin with the prescribed
images of the 24 known points.  Whenever $T(a)$ is known, the
equations force the images of $X(a)$ and $Y(a)$, respectively,
to be $(RXR^{-1})(T(a))$ and $Y(T(a))$.  A conflict of either
images or inverse images rejects the candidate.  Transitivity
of $\langle X,Y\rangle$ ensures that propagation reaches every
point.  Therefore this procedure is both necessary and
sufficient for the stated finite intertwiner problem, and it
has at most one solution.  Its exact result is sixteen accepted
values and four distinct intertwiners.  Direct multiplication
checks that the sixteen values generate a group of order
sixteen and that all fifteen nonidentity elements are
involutions.  The subgroup is therefore $C_2^4$.

Finally join the two degree-72 representations on 144 points.
Its image has order 995328, equal to both projections; hence
both projections are isomorphisms and the marked kernels
coincide.  Its derived image has order 82944.  Membership
testing of the $16^2$ pairs in this derived image gives
exactly sixteen pairs, bijective under both projections.
Reading their already unique intertwiners gives the four pairs
in the statement.  This is a same-word calculation in finite
quotients and needs no assertion that finite candidates lift
to actual gentle or full-GT elements.
\end{proof}

\paragraph{Exact certificates and scope.}
The exhaustive source-marked calculation is
\path{agents/characters/code/hurwitz_mixed72_partial.py},
with JSON and NPZ records of the same stem.  Its NPZ records the
global point labels, both restricted generators, the 24 marked
fixed points, all sixteen $R$ values, their intertwiners, and
the four lexicographically ordered distinct intertwiners.
The graph comparison is
\path{agents/characters/code/hurwitz_mixed72_joint.py}, with
\path{agents/characters/records/hurwitz_mixed72_joint_probe.json}.
The actual-source inputs are P18, P21, H4, and the supplied
FP.37.2 marked $S_4$ formula.  The finite comparison is stronger
than the compact order-576 quotient of C21, but it still leaves
both cubic states possible.

\paragraph{Further local tests: an explicitly recorded failure.}
For each of the 72 global permutations retained after C25 in
the subsequent global calculation, all other mixed-pair orbit
restrictions were tested independently.  Every restriction
normalizes its marked image and sends $X$ to a conjugate under
that image's derived subgroup.  Explicit local conjugators
were found for every case; no rejection was inferred from a
search cutoff.  The exact records are
\path{agents/characters/records/hurwitz_global_mixed_local.json}
and the accompanying NPZ, generated by the code of the same
stem.  Thus these separate local tests do not reduce the 72
global candidates.  Common-word compatibility across the
components and the full equation involving $C$ remain
additional constraints.
